\documentclass[
  man,
  floatsintext,
  longtable,
  nolmodern,
  notxfonts,
  notimes,
  colorlinks=true,linkcolor=blue,citecolor=blue,urlcolor=blue]{apa7}

\usepackage{amsmath}
\usepackage{amssymb}
\usepackage{setspace}



\RequirePackage{longtable}
\RequirePackage{threeparttablex}

\usepackage{framed}


\makeatletter
\renewcommand{\paragraph}{\@startsection{paragraph}{4}{\parindent}%
	{0\baselineskip \@plus 0.2ex \@minus 0.2ex}%
	{-.5em}%
	{\normalfont\normalsize\bfseries\typesectitle}}

\renewcommand{\subparagraph}[1]{\@startsection{subparagraph}{5}{0.5em}%
	{0\baselineskip \@plus 0.2ex \@minus 0.2ex}%
	{-\z@\relax}%
	{\normalfont\normalsize\bfseries\itshape\hspace{\parindent}{#1}\textit{\addperi}}{\relax}}
\makeatother




\usepackage{longtable, booktabs, multirow, multicol, colortbl, hhline, caption, array, float, xpatch}
\usepackage{subcaption}


\renewcommand\thesubfigure{\Alph{subfigure}}
\setcounter{topnumber}{2}
\setcounter{bottomnumber}{2}
\setcounter{totalnumber}{4}
\renewcommand{\topfraction}{0.85}
\renewcommand{\bottomfraction}{0.85}
\renewcommand{\textfraction}{0.15}
\renewcommand{\floatpagefraction}{0.7}

\usepackage{tcolorbox}
\tcbuselibrary{listings,theorems, breakable, skins}
\usepackage{fontawesome5}

\definecolor{quarto-callout-color}{HTML}{909090}
\definecolor{quarto-callout-note-color}{HTML}{0758E5}
\definecolor{quarto-callout-important-color}{HTML}{CC1914}
\definecolor{quarto-callout-warning-color}{HTML}{EB9113}
\definecolor{quarto-callout-tip-color}{HTML}{00A047}
\definecolor{quarto-callout-caution-color}{HTML}{FC5300}
\definecolor{quarto-callout-color-frame}{HTML}{ACACAC}
\definecolor{quarto-callout-note-color-frame}{HTML}{4582EC}
\definecolor{quarto-callout-important-color-frame}{HTML}{D9534F}
\definecolor{quarto-callout-warning-color-frame}{HTML}{F0AD4E}
\definecolor{quarto-callout-tip-color-frame}{HTML}{02B875}
\definecolor{quarto-callout-caution-color-frame}{HTML}{FD7E14}

%\newlength\Oldarrayrulewidth
%\newlength\Oldtabcolsep


\usepackage{hyperref}
\usepackage[protrusion=true]{microtype}



\providecommand{\tightlist}{%
  \setlength{\itemsep}{0pt}\setlength{\parskip}{0pt}}
\usepackage{longtable,booktabs,array}
\usepackage{calc} % for calculating minipage widths
% Correct order of tables after \paragraph or \subparagraph
\usepackage{etoolbox}
\makeatletter
\patchcmd\longtable{\par}{\if@noskipsec\mbox{}\fi\par}{}{}
\makeatother
% Allow footnotes in longtable head/foot
\IfFileExists{footnotehyper.sty}{\usepackage{footnotehyper}}{\usepackage{footnote}}
\makesavenoteenv{longtable}

\usepackage{graphicx}
\makeatletter
\newsavebox\pandoc@box
\newcommand*\pandocbounded[1]{% scales image to fit in text height/width
  \sbox\pandoc@box{#1}%
  \Gscale@div\@tempa{\textheight}{\dimexpr\ht\pandoc@box+\dp\pandoc@box\relax}%
  \Gscale@div\@tempb{\linewidth}{\wd\pandoc@box}%
  \ifdim\@tempb\p@<\@tempa\p@\let\@tempa\@tempb\fi% select the smaller of both
  \ifdim\@tempa\p@<\p@\scalebox{\@tempa}{\usebox\pandoc@box}%
  \else\usebox{\pandoc@box}%
  \fi%
}
% Set default figure placement to htbp
\def\fps@figure{htbp}
\makeatother


% definitions for citeproc citations
\NewDocumentCommand\citeproctext{}{}
\NewDocumentCommand\citeproc{mm}{%
  \begingroup\def\citeproctext{#2}\cite{#1}\endgroup}
\makeatletter
 % allow citations to break across lines
 \let\@cite@ofmt\@firstofone
 % avoid brackets around text for \cite:
 \def\@biblabel#1{}
 \def\@cite#1#2{{#1\if@tempswa , #2\fi}}
\makeatother
\newlength{\cslhangindent}
\setlength{\cslhangindent}{1.5em}
\newlength{\csllabelwidth}
\setlength{\csllabelwidth}{3em}
\newenvironment{CSLReferences}[2] % #1 hanging-indent, #2 entry-spacing
 {\begin{list}{}{%
  \setlength{\itemindent}{0pt}
  \setlength{\leftmargin}{0pt}
  \setlength{\parsep}{0pt}
  % turn on hanging indent if param 1 is 1
  \ifodd #1
   \setlength{\leftmargin}{\cslhangindent}
   \setlength{\itemindent}{-1\cslhangindent}
  \fi
  % set entry spacing
  \setlength{\itemsep}{#2\baselineskip}}}
 {\end{list}}
\usepackage{calc}
\newcommand{\CSLBlock}[1]{\hfill\break\parbox[t]{\linewidth}{\strut\ignorespaces#1\strut}}
\newcommand{\CSLLeftMargin}[1]{\parbox[t]{\csllabelwidth}{\strut#1\strut}}
\newcommand{\CSLRightInline}[1]{\parbox[t]{\linewidth - \csllabelwidth}{\strut#1\strut}}
\newcommand{\CSLIndent}[1]{\hspace{\cslhangindent}#1}





\usepackage{newtx}

\defaultfontfeatures{Scale=MatchLowercase}
\defaultfontfeatures[\rmfamily]{Ligatures=TeX,Scale=1}





\title{Discussion Draft}


\shorttitle{Discussion Draft}


\usepackage{etoolbox}






\author{Name}



\affiliation{
{Universiteit}}




\leftheader{Name}



\abstract{Hallo }

\keywords{open science, preregistration, preregistration
deviation, reproducibility crisis, p-hacking, simulation
study, metascience, type I error, type II error, social
science, psychology}

\authornote{\par{\addORCIDlink{Name}{0000-0000-0000-0000}} 

\par{   The author(s) declare that there were no conflicts of interest
with respect to the authorship or the publication of this
article.    Author roles were classified using the Contributor Role
Taxonomy (CRediT; https://credit.niso.org/) as
follows:  Name:   conceptualization, methodology, investigation, formal
analysis, software, visualization, writing}
\par{Correspondence concerning this article should be addressed to Name}
}

\makeatletter
\let\endoldlt\endlongtable
\def\endlongtable{
\hline
\endoldlt
}
\makeatother

\urlstyle{same}



\usepackage{booktabs}
\usepackage{longtable}
\usepackage{array}
\usepackage{multirow}
\usepackage{wrapfig}
\usepackage{float}
\usepackage{colortbl}
\usepackage{pdflscape}
\usepackage{tabu}
\usepackage{threeparttable}
\usepackage{threeparttablex}
\usepackage[normalem]{ulem}
\usepackage{makecell}
\usepackage{xcolor}
\AtBeginDocument{\thispagestyle{empty}}
\usepackage{tabularx}
\usepackage{threeparttable}
\usepackage{ragged2e}
\usepackage{booktabs}
\makeatletter
\@ifpackageloaded{caption}{}{\usepackage{caption}}
\AtBeginDocument{%
\ifdefined\contentsname
  \renewcommand*\contentsname{Table of contents}
\else
  \newcommand\contentsname{Table of contents}
\fi
\ifdefined\listfigurename
  \renewcommand*\listfigurename{List of Figures}
\else
  \newcommand\listfigurename{List of Figures}
\fi
\ifdefined\listtablename
  \renewcommand*\listtablename{List of Tables}
\else
  \newcommand\listtablename{List of Tables}
\fi
\ifdefined\figurename
  \renewcommand*\figurename{Figure}
\else
  \newcommand\figurename{Figure}
\fi
\ifdefined\tablename
  \renewcommand*\tablename{Table}
\else
  \newcommand\tablename{Table}
\fi
}
\@ifpackageloaded{float}{}{\usepackage{float}}
\floatstyle{ruled}
\@ifundefined{c@chapter}{\newfloat{codelisting}{h}{lop}}{\newfloat{codelisting}{h}{lop}[chapter]}
\floatname{codelisting}{Listing}
\newcommand*\listoflistings{\listof{codelisting}{List of Listings}}
\makeatother
\makeatletter
\makeatother
\makeatletter
\@ifpackageloaded{caption}{}{\usepackage{caption}}
\@ifpackageloaded{subcaption}{}{\usepackage{subcaption}}
\makeatother

% From https://tex.stackexchange.com/a/645996/211326
%%% apa7 doesn't want to add appendix section titles in the toc
%%% let's make it do it
\makeatletter
\xpatchcmd{\appendix}
  {\par}
  {\addcontentsline{toc}{section}{\@currentlabelname}\par}
  {}{}
\makeatother

%% Disable longtable counter
%% https://tex.stackexchange.com/a/248395/211326

\usepackage{etoolbox}

\makeatletter
\patchcmd{\LT@caption}
  {\bgroup}
  {\bgroup\global\LTpatch@captiontrue}
  {}{}
\patchcmd{\longtable}
  {\par}
  {\par\global\LTpatch@captionfalse}
  {}{}
\apptocmd{\endlongtable}
  {\ifLTpatch@caption\else\addtocounter{table}{-1}\fi}
  {}{}
\newif\ifLTpatch@caption
\makeatother

\begin{document}

\maketitle




\setlength\LTleft{0pt}



\ifdefined\Shaded\renewenvironment{Shaded}{\begin{snugshade}\begin{singlespace}}{\end{singlespace}\end{snugshade}}\fi

\setlength{\cslhangindent}{0.5in}

\section{Discussion}\label{discussion}

The present multiverse analysis aimed to quantify the relative
contribution of harmful versus helpful analytical decisions to effect
size variation for the association between marital status and
cardiovascular events, based on the many-analyst project by Kowall et
al. (\citeproc{ref-kowall_marital_2025}{2025}). Across 4,913
specifications, risk ratio estimates ranged from 0.71 to 1.68, with
79.6\% of the universes yielding significant results. Marriage was found
to be both protective and harmful for cardiovascular health depending on
the analytical specification. This underscores that different analytical
pathways can lead to substantially different conclusions --- a pattern
consistent with findings from earlier many-analyst projects
(\citeproc{ref-bastiaansen_time_2020}{Bastiaansen et al., 2020};
\citeproc{ref-botvinik-nezer_variability_2020}{Botvinik-Nezer et al.,
2020}; \citeproc{ref-breznau_observing_2022}{Breznau et al., 2022};
\citeproc{ref-gould_same_2025}{Gould et al., 2025};
\citeproc{ref-kowall_marital_2025}{Kowall et al., 2025}; e.g.,
\citeproc{ref-silberzahn_many_2018}{Silberzahn et al., 2018}).

Although the Anderson-Darling tests showed statistically significant
distributional differences for all decisions, indicating they were all
sensitive, statistical model and adjustment set emerged as the most
prominent drivers of variation in RR distributions. When applying the
AB-framework classification, both a harmful source of
variation---statistical model---and a helpful source---adjustment
set---contributed meaningfully to effect size variability. All other
decisions, despite reaching statistical significance, showed minimal
visual distributional differences across options, suggesting their
impact on effect size variability was small. This was reinforced by the
descriptive plots and T.AD values, which suggest that only statistical
model and adjustment set had meaningful impact.

\subsection{Helpful Sources of
Variation}\label{helpful-sources-of-variation}

A closer examination of adjustment set sensitivity using
Kolmogorov-Smirnov tests showed that age, sex, and not including any
confounders were the most important drivers of distributional
differences in the results. Age and sex showed opposing patterns:
controlling for age attenuated the protective effect of marriage, while
controlling for sex strengthened it. A possible explanation for this
pattern is the distribution of these variables over the groups in the
underlying data. Unmarried individuals were on average four years older
and predominantly female (69\%), meaning the married and unmarried
groups differed systematically in age and sex. Without adjustment for
these variables, the estimated effect of marriage on cardiovascular
events therefore reflects not only the effect of marriage itself, but
also these compositional differences between groups. The younger average
age of married individuals may have inflated the apparent protective
effect of marriage, while the higher proportion of males in this group
may have attenuated it, given that men have greater cardiovascular risk
(\citeproc{ref-pana_sex-specific_2024}{Pana et al., 2024}).

From a causal inference perspective, the inclusion of covariates shifts
the estimation from a total effect to a direct effect
(\citeproc{ref-rohrer_thinking_2018}{Rohrer, 2018}). Without adjustment,
the comparison between married and unmarried individuals mixes the
effect of marriage with differences in age and sex composition between
the groups, rather than isolating the effect of marriage itself.
Adjusting for age, for example, estimates the effect of marriage among
people of the same age, and similarly for sex. The observed sensitivity
to these covariates therefore reflects a genuine difference in what is
being estimated, not arbitrary noise---the kind of variation the
AB-framework would characterise as helpful.

Consistent with the T.AD values, the pattern observed for the
no-confounders condition (see fig-dens-ADJ when `none' was excluded)
closely resembled that of the age-exclusion condition, suggesting that
age was the most important variable affecting the decision sensitivity
of the adjustment set. When no confounders were included, risk ratio
estimates were lower, and when any confounder was included, the
estimates shifted towards higher values. This supports the
interpretation that variation in the results in this case represents
sensitivity to the causal structure of the data.

\subsection{Harmful Sources of
Variation}\label{harmful-sources-of-variation}

When looking at the harmful decisions, statistical models were the most
important drivers of variation in the results. Especially Cox and
log-binomial models resulted in visually different distributions of risk
ratios (see fig-dens-plot). It is important to note that when setting up
the multiverse, it was evident that the choice of statistical models was
not an independent decision. For example, only three out of seven
outcome operationalisations contained time-to-event data, and only two
out of four data types were compatible with Cox models. Choosing this
statistical model therefore implicitly constrained other facets of the
analysis. Crucially, even after conversion to risk ratios, Cox models
consistently yielded substantially higher estimates than other modelling
approaches, although these were often non-significant. This suggests
that model assumptions can even fundamentally alter the conclusion.

Additionally, 1,102 log-binomial models failed to converge---whereas
only 232 were successful. Non-convergence occurred most often when the
adjustment sets increased in size, which connects to a known limitation
of the model (\citeproc{ref-zhu_refinements_2024}{Zhu et al., 2024}).
Log-binomial models restrict predicted disease probabilities to lie
between zero and one, which in turn bounds the values the regression
coefficients can take
(\citeproc{ref-williamson_log-binomial_2013}{Williamson et al., 2013}).
Convergence failure arises when the maximum likelihood solution sits on
the boundary of the parameter space, when at least one combination of
covariates results in an estimated probability of exactly one
(\citeproc{ref-zhu_refinements_2024}{Zhu et al., 2024}). This may become
more likely as adjustment sets grow, since larger sets may be more
likely to produce covariate combinations that hit this boundary.
Notably, convergence behaviour for log-binomial models may be
software-specific
(\citeproc{ref-williamson_log-binomial_2013}{Williamson et al., 2013}).
In the present study, all analyses were conducted in R, so this was not
accounted for, meaning that some of the observed non-convergence may
reflect software choices rather than properties of the model itself.

Choosing a log-binomial model is therefore not an independent analytical
decision, as its convergence depends at least partly on size of the
adjustment set. A similar dependency arises with GEE models, which
cannot handle missing data and therefore force an explicit decision
about missingness, whereas other models default to complete-case
analysis when nothing is specified. This means that model choice in this
case again requires choices elsewhere in the analytical pathway.
Importantly, choices regarding missing data in itself can also have
implications that go beyond arbitrary noise. If data are not missing at
random, listwise deletion could result in not including a subpopulation
with a certain characteristic in the analysis. Moreover, when
confounders are missing not at random, causal effects are not
identifiable (\citeproc{ref-yang_causal_2019}{Yang et al., 2019}).

Taken together, these examples reveal a fundamental limitation of the
AB-framework: analytical decisions do not exist in isolation, so
categorising them as such is problematic. Decisions about statistical
models are entangled with data structure requirements, missing data
handling and variable operationalisation. This means that many harmful
decisions may in fact carry substantive assumptions that cannot be
separated from potential arbitrary variation the framework was designed
to identify.

\subsection{Comparison With Original Many-Analyst
Study}\label{comparison-with-original-many-analyst-study}

The range of effect sizes found in this study was broadly in line with
that of the original many-analyst project by Kowall et al.
(\citeproc{ref-kowall_marital_2025}{2025}), where effect sizes ranged
from 0.72 (OR) to 1.31 (HR). After applying the risk ratio conversions
used in the present study, the original effect sizes would correspond to
a range of approximately 0.85 to 1.21, indicating the original study
captured a slightly narrower range than the present multiverse. This is
not surprising, as the original study only included sixteen analyses,
compared to 4,913 analyses.

Despite the overlap in effect size ranges, the two studies diverged in
identifying the main sources of variability. In contrast to the current
study, Kowall et al. (\citeproc{ref-kowall_marital_2025}{2025})
identified the data type (cross-sectional versus longitudinal) as the
main source of variation, with statistical model and the adjustment set
playing only a minor role---the opposite of the present findings. This
inconsistency could reflect a fundamental difference between the two
approaches. Whereas Kowall et al.
(\citeproc{ref-kowall_marital_2025}{2025}) varied one decision at a time
under a fixed configuration, this multiverse analysis combined all
decisions simultaneously and captured their combined influence across
the entire decision space.

Another potential explanation for the discrepancy between studies is the
finding that adjustment set sensitivity was driven primarily by age and
sex. In Kowall et al. (\citeproc{ref-kowall_marital_2025}{2025}), twelve
out of sixteen researchers included both age and sex in their analyses,
and the additional sensitivity analyses also held these variables
constant. Most researchers in the original study thus estimated the
direct effect of marriage on cardiovascular events, leaving little data
available about the effect of not including these confounders. This may
potentially have attenuated the observed sensitivity to adjustment set
decisions.

Rather than contradicting the original study, the present findings
extend it. When analytical decisions are varied jointly, rather than in
isolation, this can provide a different assessment of decision
sensitivity.

\subsection{Limitations}\label{limitations}

A first limitation concerns the inferential method used to assess
decision sensitivity. The k-sample Anderson-Darling test proved to be
sensitive to sample size . With 4,913 specifications, even minimal
distributional differences resulted in significant results. The
universal significance across all decisions therefore likely reflects
the number of specifications rather than substantive distributional
differences. While the AD test yielded a useful measure of
between-decision sensitivity, combining it with descriptive methods
rather than relying on p-values alone would provide a more informative
picture of decision sensitivity. Moreover, future research should
explore inferential methods that are less susceptible to the large
sample sizes characteristic of multiverse analyses, and explicitly
account for the dependency structure between decisions and
specifications.

A second limitation is the generalisability of the results of the
multiverse study. The relative contribution of harmful versus helpful
analytical decisions will likely vary across studies, given the wide
variety of topics and research questions across many-analyst projects
(\citeproc{ref-gould_same_2025}{Gould et al., 2025};
\citeproc{ref-hoogeveen_many-analysts_2023}{Hoogeveen et al., 2023};
\citeproc{ref-kowall_marital_2025}{Kowall et al., 2025}; e.g.,
\citeproc{ref-silberzahn_many_2018}{Silberzahn et al., 2018}). It does
not seem realistic to expect that the results of this study directly
apply to other projects. The decisions that proved to be most
influential here---statistical model and adjustment set---are not
necessarily the main sources of variation in other contexts, or even in
the same context when decisions are varied one at a time rather than
jointly, as the comparison with Kowall et al.
(\citeproc{ref-kowall_marital_2025}{2025}) illustrates. Replicating this
multiverse analysis across different many-analyst projects, preferably
with a wide range of research questions and data structures, might
reveal more general patterns in decision sensitivity. However, assuming
that such patterns exist may perhaps already be too bold.

Another limitation is the conversion of effect sizes. While this has
benefits for comparability across statistical models, the converged risk
ratios are always an approximation. According to VanderWeele
(\citeproc{ref-vanderweele_optimal_2020}{2020}), the square root
transformation of the odds ratio gives a reasonably good approximation
and differs from the true risk ratio by at most 25\%, when the
prevalence lies within 20\%--80\%. In case of hazard ratios, the
conversion will place the value within 16\% of the true risk ratios when
the prevalence lies within 20\%--80\%, and 45\% of the risk ratio if the
prevalence is between 10\%--90\%, which is the case in the present
study. This highlights that even though converting effect sizes can have
benefits for inter-model comparisons, at the same time this may have
introduced, or reduced, arbitrary variance in the results.

\subsection{Future Recommendations}\label{future-recommendations}

The current findings raise a broader question: what should analysts do
when analytical choices can fundamentally alter the conclusions of their
research? Multiverse analysis cannot determine which specification is
``correct'', and the present study showed that substantial variability
in effect estimates was found even within multiple defensible analytical
paths. To ensure that the effects of analytical variability can be
assessed properly, it is recommended to perform multiverse analysis, not
only as a measure of robustness, but also as a transparent
representation of the space of plausible conclusions
(\citeproc{ref-botvinik-nezer_variability_2020}{Botvinik-Nezer et al.,
2020}).

Future research should explore inferential methods to quantify decision
sensitivity that are more suitable and robust to the large samples in
multiverse analysis. Additionally, extending these methods to allow the
derivation of a unique result that integrates all uncertain decisions
remains an open area of research
(\citeproc{ref-aczel_consensus-based_2021}{Aczel et al., 2021}). In
addition to developing new methods, researchers should emphasize ground
their decisions in theory and causal reasoning rather than idiosyncratic
choices. Notably, only three of the sixteen teams in Kowall et al.
(\citeproc{ref-kowall_marital_2025}{2025}) used directed acyclic graphs
(DAGs), and even those differed significantly. Specifying analytical
decisions in pre-registrations, constructing DAGs before starting the
analysis, and making causal assumptions explicit could help reduce
researchers' unknown degrees of freedom
(\citeproc{ref-botvinik-nezer_variability_2020}{Botvinik-Nezer et al.,
2020}). In combination with the openly sharing data and code, this would
not eliminate analytical variability, but would ensure that its effects
can be assessed in a transparent manner to meaningfully interpret
variation in the results and associated substantive conclusions.

\newpage

\section{References}\label{references}

\phantomsection\label{refs}
\begin{CSLReferences}{1}{0}
\bibitem[\citeproctext]{ref-aczel_consensus-based_2021}
Aczel, B., Szaszi, B., Nilsonne, G., Van Den Akker, O. R., Albers, C.
J., Van Assen, M. A., Bastiaansen, J. A., Benjamin, D., Boehm, U.,
Botvinik-Nezer, R., Bringmann, L. F., Busch, N. A., Caruyer, E.,
Cataldo, A. M., Cowan, N., Delios, A., Van Dongen, N. N., Donkin, C.,
Van Doorn, J. B., \ldots{} Wagenmakers, E.-J. (2021). Consensus-based
guidance for conducting and reporting multi-analyst studies.
\emph{eLife}, \emph{10}, e72185.
\url{https://doi.org/10.7554/eLife.72185}

\bibitem[\citeproctext]{ref-bastiaansen_time_2020}
Bastiaansen, J. A., Kunkels, Y. K., Blaauw, F. J., Boker, S. M.,
Ceulemans, E., Chen, M., Chow, S.-M., De Jonge, P., Emerencia, A. C.,
Epskamp, S., Fisher, A. J., Hamaker, E. L., Kuppens, P., Lutz, W.,
Meyer, M. J., Moulder, R., Oravecz, Z., Riese, H., Rubel, J., \ldots{}
Bringmann, L. F. (2020). Time to get personal? {The} impact of
researchers choices on the selection of treatment targets using the
experience sampling methodology. \emph{Journal of Psychosomatic
Research}, \emph{137}, 110211.
\url{https://doi.org/10.1016/j.jpsychores.2020.110211}

\bibitem[\citeproctext]{ref-botvinik-nezer_variability_2020}
Botvinik-Nezer, R., Holzmeister, F., Camerer, C. F., Dreber, A., Huber,
J., Johannesson, M., Kirchler, M., Iwanir, R., Mumford, J. A., Adcock,
R. A., Avesani, P., Baczkowski, B. M., Bajracharya, A., Bakst, L., Ball,
S., Barilari, M., Bault, N., Beaton, D., Beitner, J., \ldots{}
Schonberg, T. (2020). Variability in the analysis of a single
neuroimaging dataset by many teams. \emph{Nature}, \emph{582}(7810),
84--88. \url{https://doi.org/10.1038/s41586-020-2314-9}

\bibitem[\citeproctext]{ref-breznau_observing_2022}
Breznau, N., Rinke, E. M., Wuttke, A., Nguyen, H. H. V., Adem, M.,
Adriaans, J., Alvarez-Benjumea, A., Andersen, H. K., Auer, D., Azevedo,
F., Bahnsen, O., Balzer, D., Bauer, G., Bauer, P. C., Baumann, M.,
Baute, S., Benoit, V., Bernauer, J., Berning, C., \ldots{} Żółtak, T.
(2022). Observing many researchers using the same data and hypothesis
reveals a hidden universe of uncertainty. \emph{Proceedings of the
National Academy of Sciences}, \emph{119}(44), e2203150119.
\url{https://doi.org/10.1073/pnas.2203150119}

\bibitem[\citeproctext]{ref-gould_same_2025}
Gould, E., Fraser, H. S., Parker, T. H., Nakagawa, S., Griffith, S. C.,
Vesk, P. A., Fidler, F., Hamilton, D. G., Abbey-Lee, R. N., Abbott, J.
K., Aguirre, L. A., Alcaraz, C., Aloni, I., Altschul, D., Arekar, K.,
Atkins, J. W., Atkinson, J., Baker, C. M., Barrett, M., \ldots{}
Zitomer, R. A. (2025). Same data, different analysts: Variation in
effect sizes due to analytical decisions in ecology and evolutionary
biology. \emph{BMC Biology}, \emph{23}(1), 35.
\url{https://doi.org/10.1186/s12915-024-02101-x}

\bibitem[\citeproctext]{ref-hoogeveen_many-analysts_2023}
Hoogeveen, S., Sarafoglou, A., Aczel, B., Aditya, Y., Alayan, A. J.,
Allen, P. J., Altay, S., Alzahawi, S., Amir, Y., Anthony, F.-V., Kwame
Appiah, O., Atkinson, Q. D., Baimel, A., Balkaya-Ince, M., Balsamo, M.,
Banker, S., Bartoš, F., Becerra, M., Beffara, B., \ldots{} Wagenmakers,
E.-J. (2023). A many-analysts approach to the relation between
religiosity and well-being. \emph{Religion, Brain \& Behavior},
\emph{13}(3), 237--283.
\url{https://doi.org/10.1080/2153599X.2022.2070255}

\bibitem[\citeproctext]{ref-kowall_marital_2025}
Kowall, B., Ahrenfeldt, L. J., Basten, J., Becher, H., Brand, T., Braun,
J., Casjens, S., Claessen, H., Denz, R., Diebner, H. H., Diexer, S.,
Eisemann, N., Furrer, E., Galetzka, W., Girschik, C., Karch, A.,
Mikolajczyk, R., Peters, M., Rospleszcz, S., \ldots{} Rübsamen, N.
(2025). Marital status and risk of cardiovascular disease -- a
multi-analyst study in epidemiology. \emph{European Journal of
Epidemiology}, \emph{40}(5), 497--509.
\url{https://doi.org/10.1007/s10654-025-01235-8}

\bibitem[\citeproctext]{ref-pana_sex-specific_2024}
Pana, T. A., Mamas, M. A., Wareham, N. J., Khaw, K.-T., Dawson, D. K.,
\& Myint, P. K. (2024). Sex-specific lifetime risk of cardiovascular
events: The {European} {Prospective} {Investigation} into
{Cancer}-{Norfolk} prospective population cohort study. \emph{European
Journal of Preventive Cardiology}, \emph{31}(2), 230--241.
\url{https://doi.org/10.1093/eurjpc/zwad283}

\bibitem[\citeproctext]{ref-rohrer_thinking_2018}
Rohrer, J. M. (2018). Thinking {Clearly} {About} {Correlations} and
{Causation}: {Graphical} {Causal} {Models} for {Observational} {Data}.
\emph{Advances in Methods and Practices in Psychological Science},
\emph{1}(1), 27--42. \url{https://doi.org/10.1177/2515245917745629}

\bibitem[\citeproctext]{ref-silberzahn_many_2018}
Silberzahn, R., Uhlmann, E. L., Martin, D. P., Anselmi, P., Aust, F.,
Awtrey, E., Bahník, Š., Bai, F., Bannard, C., Bonnier, E., Carlsson, R.,
Cheung, F., Christensen, G., Clay, R., Craig, M. A., Dalla Rosa, A.,
Dam, L., Evans, M. H., Flores Cervantes, I., \ldots{} Nosek, B. A.
(2018). Many {Analysts}, {One} {Data} {Set}: {Making} {Transparent}
{How} {Variations} in {Analytic} {Choices} {Affect} {Results}.
\emph{Advances in Methods and Practices in Psychological Science},
\emph{1}(3), 337--356. \url{https://doi.org/10.1177/2515245917747646}

\bibitem[\citeproctext]{ref-vanderweele_optimal_2020}
VanderWeele, T. J. (2020). Optimal approximate conversions of odds
ratios and hazard ratios to risk ratios. \emph{Biometrics},
\emph{76}(3), 746--752. \url{https://doi.org/10.1111/biom.13197}

\bibitem[\citeproctext]{ref-williamson_log-binomial_2013}
Williamson, T., Eliasziw, M., \& Fick, G. H. (2013). Log-binomial
models: Exploring failed convergence. \emph{Emerging Themes in
Epidemiology}, \emph{10}, 14.
\url{https://doi.org/10.1186/1742-7622-10-14}

\bibitem[\citeproctext]{ref-yang_causal_2019}
Yang, S., Wang, L., \& Ding, P. (2019). Causal inference with
confounders missing not at random. \emph{Biometrika}, \emph{106}(4),
875--888. \url{https://doi.org/10.1093/biomet/asz048}

\bibitem[\citeproctext]{ref-zhu_refinements_2024}
Zhu, C., Hosmer, D. W., Stankovich, J., Wills, K., \& Blizzard, L.
(2024). Refinements on the exact method to solve the numerical
difficulties in fitting the log binomial regression model for estimating
relative risk. \emph{Communications in Statistics - Theory and Methods},
\emph{53}(23), 8359--8375.
\url{https://doi.org/10.1080/03610926.2023.2284674}

\end{CSLReferences}

\begin{verbatim}
# A tibble: 89 x 7
   Decision          Option        N_valid N_failed `Mean RR` `RR Range` `% Sig`
   <chr>             <chr>           <int>    <int>     <dbl> <chr>      <chr>  
 1 Statistical Model Log-binomial      232     1102      0.83 [0.71, 0.~ 100%   
 2 Statistical Model Mixed effect~    1278        0      0.92 [0.73, 1.~ 77%    
 3 Statistical Model Discrete-tim~     832        0      0.92 [0.76, 1.~ 74%    
 4 Statistical Model GEE              1052      221      0.95 [0.8, 1.0~ 85%    
 5 Statistical Model Logistic reg~    1188      164      0.95 [0.8, 1.1~ 93%    
 6 Statistical Model Cox PH            331        0      1.08 [0.84, 1.~ 24%    
 7 Data Type         Cross-sectio~    1237      403      0.93 [0.74, 1.~ 91%    
 8 Data Type         Longitudinal~    1669      563      0.93 [0.71, 1.~ 69%    
 9 Data Type         Longitudinal~    1869      411      0.95 [0.75, 1.~ 82%    
10 Data Type         Cross-sectio~     138      110      0.97 [0.77, 1.~ 86%    
# i 79 more rows
\end{verbatim}






\end{document}
